\documentclass[11pt]{article}

% ── Packages ────────────────────────────────────────────────────────────────
\usepackage[margin=1in]{geometry}
\usepackage{amsmath, amssymb, amsthm}
\usepackage{mathtools}
\usepackage{enumitem}
\usepackage{hyperref}
\usepackage{parskip}

% ── Theorem environments ─────────────────────────────────────────────────────
\newtheorem{theorem}{Theorem}[section]
\newtheorem{lemma}[theorem]{Lemma}
\newtheorem{corollary}[theorem]{Corollary}
\newtheorem{proposition}[theorem]{Proposition}
\theoremstyle{definition}
\newtheorem{definition}[theorem]{Definition}
\newtheorem{example}[theorem]{Example}
\theoremstyle{remark}
\newtheorem{remark}[theorem]{Remark}

% ── Custom commands ───────────────────────────────────────────────────────────
\newcommand{\R}{\mathbb{R}}
\newcommand{\E}{\mathbb{E}}
\newcommand{\norm}[1]{\left\lVert#1\right\rVert}
\newcommand{\inner}[2]{\langle #1,\, #2 \rangle}
\newcommand{\T}{^{\top}}
\newcommand{\K}{K}  % number of classes

% ── Document ─────────────────────────────────────────────────────────────────
\begin{document}

\title{\textbf{Multiclassification: Decision}\\[0.4em]
       \large Notes Template}
\author{}
\date{}
\maketitle
\tableofcontents
\newpage

% ─────────────────────────────────────────────────────────────────────────────
\section{Problem}
% ─────────────────────────────────────────────────────────────────────────────

\subsection{Setup}

% TODO: Describe the multiclass decision problem here.
% Example prompts:
%   - What is the input/output space?
%   - How many classes K are there?
%   - What is the loss function (0-1 loss, cross-entropy, …)?
%   - What decision rule maps class posteriors to a predicted label?

\begin{definition}[Multiclass Decision Problem]
    % TODO: Fill in the formal definition.
    Let $\mathcal{X} \subseteq \R^{d}$ be the input space and
    $\mathcal{Y} = \{1, 2, \ldots, \K\}$ the label space with $\K \geq 2$ classes.
    Given a probabilistic model $p(y \mid x)$,
    the decision problem is to choose a label $\hat{y} \in \mathcal{Y}$ such that \ldots
\end{definition}

\subsection{Assumptions}

\begin{enumerate}[label=\textbf{A\arabic*.}]
    \item % TODO: State assumption 1 (e.g., access to class posteriors).
    \item % TODO: State assumption 2 (e.g., loss structure).
    \item % TODO: Add further assumptions as needed.
\end{enumerate}

\subsection{Objective}

% TODO: State the decision objective (e.g., minimise expected loss / Bayes risk).
\[
    % TODO: Replace with the actual decision rule.
    \hat{y} \;=\; \arg\max_{k \in \mathcal{Y}} \; p(y = k \mid x)
\]

% ─────────────────────────────────────────────────────────────────────────────
\section{Derivation}
% ─────────────────────────────────────────────────────────────────────────────

The following derivation establishes the optimal multiclass decision rule
step by step.

% ── Step 1 ────────────────────────────────────────────────────────────────────
\subsection{Step 1: [Title]}

% TODO: Describe what this step accomplishes (e.g., define the Bayes risk).

\begin{proof}
    % TODO: Fill in the proof for Step 1.
    \[
        \text{(derivation goes here)}
    \]
\end{proof}

% ── Step 2 ────────────────────────────────────────────────────────────────────
\subsection{Step 2: [Title]}

% TODO: Describe what this step accomplishes (e.g., expand the risk over classes).

\begin{proof}
    % TODO: Fill in the proof for Step 2.
    \begin{align}
        R(\hat{y}) &= \E_{x}\!\left[\sum_{k=1}^{\K} \mathcal{L}(\hat{y}, k)\, p(y=k \mid x)\right] \notag \\
                   &= \text{(simplify)} \label{eq:step2}
    \end{align}
\end{proof}

% ── Step 3 ────────────────────────────────────────────────────────────────────
\subsection{Step 3: [Title]}

% TODO: Describe what this step accomplishes (e.g., minimise pointwise over x).

\begin{proof}
    % TODO: Fill in the proof for Step 3.
    \[
        \text{(derivation goes here)}
    \]
\end{proof}

% ── Final Result ──────────────────────────────────────────────────────────────
\subsection{Final Result}

\begin{theorem}[Optimal Multiclass Decision Rule]
    % TODO: State the final theorem.
    Under assumptions \textbf{A1}--\textbf{A3} and a 0-1 loss,
    the Bayes-optimal decision rule is
    \[
        % TODO: Replace with the actual result.
        \hat{y}^{*}(x) \;=\; \arg\max_{k \in \{1,\ldots,\K\}} \; p(y = k \mid x).
    \]
\end{theorem}

\begin{proof}
    Follows directly from Steps 1--3. \qed
\end{proof}

% ─────────────────────────────────────────────────────────────────────────────
\end{document}
